% Part: first-order-logic
% Chapter: introduction
% Section: formulas

\documentclass[../../../include/open-logic-section]{subfiles}

\begin{document}

\olfileid{fol}{int}{fml}

\section{\usetoken{P}{formula}}

رویکردی که برای تعیین دقیق !!{sentence}s منطق مرتبهٔ اول و حل مسائل
ناشی از کاربرد !!{variable}s به کار خواهیم برد چنین است. نخست مجموعهٔ
\emph{متفاوتی} از عبارت‌ها، یعنی !!{formula}s، را تعریف می‌کنیم. پس از
آن می‌توانیم نقشی را که !!{variable}s در آن‌ها دارند بررسی کنیم و بر
پایهٔ چند ایدهٔ دیگر، یعنی !!{variable}s «آزاد» و «مقید»، تعریف کنیم
!!a{sentence} چیست (یعنی !!a{formula} بدون !!{variable}s آزاد). این
کار را نه‌فقط از آن رو انجام می‌دهیم که تعریف «!!{sentence}» را
آسان‌تر می‌کند، بلکه از آن رو که برای شیوهٔ تعریف مفهوم معنایی ارضا
ضروری خواهد بود.

«!!{formula}» را برای یک زبان سادهٔ مرتبهٔ اول تعریف کنیم؛ زبانی که
فقط شامل یک !!{predicate}~$\Obj P$، یک !!{constant}~$\Obj a$ و نمادهای
منطقی $\lnot$، $\land$ و~$\lexists$ است. تعریف‌های کامل ما بسیار کلی‌تر
خواهند بود: بی‌نهایت !!{predicate} و !!{constant} را مجاز خواهیم دانست.
در واقع، !!{function}s را نیز بررسی خواهیم کرد که می‌توان آن‌ها را با
!!{constant}s و !!{variable}s ترکیب کرد تا «ترم» ساخت. فعلاً
$\Obj a$ و متغیرها تنها ترم‌های ما خواهند بود. البته به بی‌نهایت
!!{variable} نیاز داریم. رسماً نمادهای $\Obj v_0$، $\Obj v_1$،
\dots، را به‌عنوان متغیر به کار خواهیم برد.

\begin{defn}
مجموعهٔ \emph{!!{formula}s}~$\Frm$ به صورت زیر تعریف می‌شود:
\begin{enumerate}
\item\ollabel{fmls-atom} $\Atom{\Obj P}{\Obj a}$ و $\Atom{\Obj
  P}{\Obj v_i}$ !!{formula} هستند ($i \in \Nat$).

\tagitem{prvNot}{\ollabel{fmls-not}اگر $!A$ !!a{formula} باشد، $\lnot !A$
  !!a{formula} است.}{}

\tagitem{prvAnd}{اگر $!A$ و $!B$ !!{formula} باشند، $(!A \land
  !B)$ !!a{formula} است.}{}

\tagitem{prvEx}{\ollabel{fmls-ex}اگر $!A$ !!a{formula} و $x$ !!a{variable}
  باشد، $\lexists[x][!A]$ !!a{formula} است.}{}

\tagitem{limitClause}{\ollabel{fmls-limit}هیچ چیز دیگری !!a{formula}
  نیست.}{}
\end{enumerate}
\end{defn}

\olref{fmls-atom} به ما می‌گوید $\Atom{\Obj P}{\Obj a}$ و
$\Atom{\Obj P}{\Obj v_i}$ برای هر $i \in \Nat$ !!{formula} هستند.
این‌ها همان !!{formula}s موسوم به \emph{اتمی} هستند و نقطهٔ آغازی به
دست می‌دهند. بندهای دیگر راه‌هایی برای ساختن !!{formula}s جدید از
فرمول‌هایی می‌دهند که پیش‌تر ساخته‌ایم. برای نمونه، بنا بر
\olref{fmls-not}، $\lnot \Atom{\Obj P}{\Obj v_2}$ !!a{formula} است،
زیرا $\Atom{\Obj P}{\Obj v_2}$ بنا بر \olref{fmls-atom} از پیش
!!a{formula} است. سپس بنا بر \olref{fmls-ex}،
$\lexists[\Obj v_2][\lnot \Atom{\Obj P}{\Obj v_2}]$ نیز
!!{formula} است، و همین‌طور ادامه می‌دهیم. \olref{fmls-limit} می‌گوید
\emph{فقط} رشته‌هایی که به این شیوه می‌سازیم !!{formula} به شمار
می‌آیند. به‌ویژه، $\lexists[\Obj v_0][\Atom{\Obj P}{\Obj a}]$ و
$\lexists[\Obj v_0][\lexists[\Obj v_0][\Atom{\Obj P}{\Obj a}]]$
\emph{فرمول به شمار می‌آیند}، اما $(\lnot \Atom{\Obj P}{\Obj a})$
به‌سبب پرانتزهای بیرونیِ زائد فرمول نیست.

این شیوهٔ تعریف !!{formula}s را \emph{تعریف استقرایی} می‌نامند و به ما
اجازه می‌دهد با نوعی برهان استقرایی به نام \emph{استقرای ساختاری}
مطالبی را دربارهٔ !!{formula}s اثبات کنیم. این مفاهیم به‌طور کلی در
\olref[mth][ind][idf]{sec} و \olref[mth][ind][sti]{sec} بررسی شده‌اند
و بهتر است پیش از ورود به برهان‌های بعدی آن‌ها را مرور کنید. ایدهٔ
اصلی این است که برای اثبات مطلبی دربارهٔ همهٔ !!{formula}s، نخست نشان
می‌دهید آن مطلب دربارهٔ !!{formula}s اتمی صادق است؛ سپس نشان می‌دهید
\emph{اگر} دربارهٔ هر !!{formula}~$!A$ (و~$!B$) صادق باشد، دربارهٔ
$\lnot !A$، $(!A \land !B)$ و $\lexists[x][!A]$ \emph{نیز} صادق است.
برای نمونه، این روش اثبات می‌کند که مطلب دربارهٔ
$\lexists[\Obj v_2][\lnot \Atom{\Obj P}{\Obj v_2}]$ صادق است: از بخش
نخست می‌دانید دربارهٔ !!{formula} اتمی~$\Atom{\Obj P}{\Obj v_2}$ صادق
است. سپس از بخش دوم به دست می‌آورید که دربارهٔ
$\lnot \Atom{\Obj P}{\Obj v_2}$ صادق است، و با کاربرد دوبارهٔ آن
می‌بینید دربارهٔ خود $\lexists[\Obj v_2][\lnot \Atom{\Obj P}{\Obj
v_2}]$ نیز صادق است. چون همهٔ !!{formula}s به‌طور استقرایی از
!!{formula}s اتمی ساخته می‌شوند، این روش برای هر فرمولی کار می‌کند.

\end{document}
